\section{Task, Data and Evaluation}
\label{sec:task}

\subsection{The task}

Each AutoJudge dataset holds the topics of a host track and the anonymized runs
submitted to it. A run contains one response per topic. A response is a list of
sentences, each with optional citations to documents. It also carries metadata that
names the run and an anonymized team. A judge reads all of this and writes a
leaderboard with one score per run, topic and measure, plus an aggregate row for
each run. A judge may emit several measures.

A complete submission has two parts. The \emph{data submission} is the set of
leaderboards the judge produces on the test sets. The \emph{code submission} is the
judge itself, packaged as a container that TIRA runs. This lets the organizers rerun
it with an LLM of their choice. An LLM judge must also reproduce its own
leaderboard with the LLM endpoint (the server that answers its calls) switched
off. It does so by replaying every call from its prompt cache, a stored record of
each prompt and the reply it received.

\subsection{Data}

Table~\ref{tab:data} lists the datasets. One test set comes from the TREC 2026
RAG (retrieval-augmented generation) track, whose responses are short cited
answers. The other comes from the TREC 2026 RAGTIME track, whose responses are
long cited reports. Teams often submit several close variants of one pipeline,
so the table counts teams as well as runs. Both test releases are in English.
They carry the text of cited documents inline, for 99.1\% and 99.7\% of
responses.

The three pilot datasets come from TREC 2025 and have human assessments. No release
ships them, though. The truth stays on the TIRA server. The only way to measure a
judge is to run it on a pilot, upload the leaderboard and read back the correlation
the organizers compute. Every pilot number in this paper was measured this way.
Of the pilots, rag25-gen is closest to rag26 and ragtime25 is closest to ragtime26.
Each topic in dragun-repgen is a news article. The report should help a reader
assess it. DRAGUN topics do not fit the AutoJudge request format, so the organizers
converted them. Next to the original article topics, the release ships a converted
file. That is the file its own configuration names. In the converted file each
topic has a question as title, a summary of the article's claims (about 1.4\,KB) as
problem statement and a fact-checker background. The copy on TIRA instead gives
the judge only the article's headline (Section~\ref{sec:dragun}).

\begin{table}
  \caption{Datasets. The test sets have no released truth. The pilots have human
  truth, held on TIRA.}
  \label{tab:data}
  \small
  \setlength{\tabcolsep}{3pt}
  \begin{tabular}{@{}llrrrr@{}}
    \toprule
    Dataset & Kind & Runs & Teams & Topics & Responses\\
    \midrule
    rag26 (test)          & short answers & 83 & 25 & 119 & 9,877\\
    ragtime26 (test)      & long reports  & 49 & 10 & 103 & 5,047\\
    \midrule
    rag25-gen (pilot)     & short answers & 51 & 18 & 105 & 5,332\\
    ragtime25 (pilot)     & long reports  & 59 & 13 & 122 & 7,198\\
    dragun-repgen (pilot) & news reports  & 28 & 9  & 30  & 840\\
    \bottomrule
  \end{tabular}
\end{table}

\subsection{How judges are scored}
\label{sec:scoring}

TIRA correlates each measure a judge emits with each human truth measure, over the
systems' aggregate scores. It uses two rank correlations. Kendall's
$\tau$~\cite{kendall1938new} measures how often two rankings put a pair of
systems in the same order. \tauap\ is the AP rank correlation~\cite{yilmaz2008new} in its
tie-aware form~\cite{urbano2017ties}. It gives more weight to agreement at the top
of the ranking. For each truth measure TIRA reports the maximum over
the judge's measures.

On 2026-10-09 the organizers re-scored every entry on rag25-gen and ragtime25
against a revised set of truth measures. For rag25-gen there are now two, Strict
Vital Score and Sub Coverage. For ragtime25 there are Nugget Coverage Macro and F1 Macro, shortened below to Nugget Coverage and F1. We leave out its third measure, Citation Support, which currently reads
0.000 for every entry. On dragun-repgen TIRA still reports a single maximum over
all its truth measures. The organizers have not described the revised measures in
detail. Where we read meaning into them below, we go by their names. We made our
design choices under the earlier scoring, which reported one maximum over all truth
measures on every pilot. This paper gives the current numbers and notes where the
two scorings disagree.

Three features of the scoring shaped both judges. First, only the best measure
counts, so emitting several well-motivated measures costs nothing. Each judge
therefore emits four or five. The same rule flatters a system with many
measures, because the best one is picked after the fact. Two of our rows in Table~\ref{tab:pilots} emit a single measure and so get no such help. These are ``rank only'' and the ensemble (Section~\ref{sec:results}).
Second, the earlier scoring also reported the minimum over all pairs of measures.
We therefore oriented every measure so that higher is better. Third, the
organizers' evaluator breaks ties in a judge's scores by input order. A coarse
score, one with few distinct values, therefore loses agreement at random. That is
why all our measures are continuous.

\subsection{Data policy}

The track limits what developers may look at, because the test responses are
themselves the test set. Section~\ref{sec:disclosure} explains how we kept to these limits. It also describes the one engineering change that followed a failure count on the test data.
